\section{总结与学习体会}

本研究参照从零构建语言模型的实现路线，完成字符处理、因果注意力、Pre-LN GPT和约3000万正文词元的随机初始化预训练。受控对照显示，预训练在相同监督预算下改善了对联任务的预测和生成指标；针对直接微调后的对应错误，最终模型引入条件词汇表示、位置连接及双向评分，并通过规则约束提高形式可靠性。赏析实验则展示了相同底座向解释性输出的任务适配。

实践中最重要的认识是，训练损失、任务表现与规则保证分别回答不同问题。嵌入和主干通过损失学习语言分布，配对监督改变生成条件，外部表示补充信息，解码规则限定候选。将这些作用分别落实并检验，比仅观察损失下降更有助于理解模型的实际能力。数据分组、监督对齐和因果核验，也使实现与结果建立了可追溯的联系。
